% Generated draft: requires mathematical and source review.
\documentclass[11pt]{article}
\usepackage[margin=25mm]{geometry}
\usepackage{fontspec}
\setmainfont{Latin Modern Roman}
\usepackage{amsmath,amssymb,mathtools,unicode-math}
\setmathfont{Latin Modern Math}
\usepackage{enumitem,needspace,xurl,hyperref}
\setlength{\parindent}{0pt}
\setlength{\parskip}{5pt}
\setlength{\emergencystretch}{3em}
\newcommand{\N}{\mathbb{N}}
\newcommand{\Z}{\mathbb{Z}}
\newcommand{\Q}{\mathbb{Q}}
\newcommand{\R}{\mathbb{R}}
\newcommand{\C}{\mathbb{C}}
\newcommand{\F}{\mathbb{F}}
\newcommand{\A}{\mathbb{A}}
\newcommand{\PP}{\mathbb{P}}
\newcommand{\E}{\mathbb{E}}
\newcommand{\eps}{\varepsilon}
\newcommand{\dd}{\,\mathrm{d}}
\begin{document}
\section*{UnsolvedMath ID: 30006018}
\textbf{Source problem code:} catalogue entry
\subsection{Definitions and mathematical statement}
Let \(g\geq 2\), and let \(\mu=(m_1,\ldots,m_n)\) be a nonempty positive integral partition of \(2g-2\), so that \(m_i\in\Z_{>0}\) and \(\sum_{i=1}^n m_i=2g-2\). Let \(C\) be a smooth connected projective complex curve of genus \(g\), and let \(z_1,\ldots,z_n\) be distinct marked points. With the zeros labeled, the stratum of holomorphic differentials with signature \(\mu\) consists of isomorphism classes of tuples
\[
(C,\omega,z_1,\ldots,z_n),\qquad
0\neq\omega\in H^0(C,K_C),\qquad
\operatorname{div}(\omega)=\sum_{i=1}^n m_i z_i.
\]
Denote this stratum by \(\mathcal H(\mu)\). Its quotient by the scaling action \(\omega\mapsto\lambda\omega\), for \(\lambda\in\C^*\), is the projectivized stratum \(\mathbb P\mathcal H(\mu)\). It parametrizes the corresponding pointed curves for which \(\sum_i m_i z_i\) is a canonical divisor. This distinguishes the differential stratum from its projectivization; the supplied extraction uses \(\mathcal H(\mu)\) for both descriptions and has corrupted scaling notation.

\textbf{Recoverable proposal, with unresolved construction.} The supplied passage proposes decomposing the stratum according to isomorphism classes of singularities associated with these pointed curves. It describes a \(\mu\)-weighted blowup at the marked points in a central fiber \(C_0\) of an isotrivial family written as \(C\times C\), followed by contraction of the proper transform of \(C_0\). The passage describes the resulting \(s\) as an isolated Gorenstein curve singularity with \(n\) rational branches and a \(\C^*\)-action. These are supplied descriptions, not conclusions established here. The family base, blowup weights, contraction, and the precise meaning of the resulting curve singularity are insufficiently specified to define a map assigning \(s\) to each pointed curve. Consequently, even the proposed decomposition is conditional on recovering that construction and its intended isomorphism relation.

For the one-point notation \((C,z)\) appearing in the passage, the Weierstrass semigroup can be defined as
\[
H(C,z)=\{0\}\cup\{k\in\Z_{>0}:\text{there exists a meromorphic function on }C\text{ whose pole divisor is }kz\}.
\]
The passage attributes to Pinkham a relationship between monomial curve singularities and loci with prescribed Weierstrass semigroup determined by monomial parametrization exponents. It does not supply the hypotheses, exponents, or a precise theorem. It also suggests that certain deformation spaces might illuminate the topology of \(\mathcal H(\mu)\), with a connection to the Kontsevich--Zorich \(K(\pi,1)\) conjecture, but specifies neither the deformation spaces nor the intended version of that conjecture.

\textbf{Statement limitation.} The supplied verification field labels this extraction \emph{unrecoverable}: it joins material from one contribution to the title and opening sentence of the next. Its square-tiled-surface material announces results and conjectures about cylinder decompositions in large genus but states no concrete question, quantifiers, asymptotic regime, or conjectured conclusion. Thus permanent ID \(30006018\) does not determine a standalone mathematical problem. The proposal above records the recoverable content without promoting it to a precise conjecture or replacing it with questions from other records.
\subsection{Short English statement}
This extraction mixes a proposed decomposition of differential strata by associated singularities with the opening of a contribution about square-tiled surfaces. The supplied verification marks it as unrecoverable. No standalone conjecture can be stated precisely without recovering the original construction and the intended question.
\subsection{Sources}
\begin{itemize}
\item Local UnsolvedMath registry, supplied complete record: permanent ID \(30006018\), exact source problem code \texttt{catalogue entry}, title \emph{Mixed Square-Tiled-Surface Talk Transition}, in the Oberwolfach Reports Open Problems collection. The record lists proposed year 2024 and supplies no proposer attribution.
\item Supplied report citation: \emph{Riemann Surfaces: Random, Flat, and Hyperbolic Geometry} (2025), Oberwolfach Reports, pp.~2369--2424, DOI \texttt{10.4171/owr/2024/41}. The source URL and the literature-source entry labeled ``OWR report'' both give \url{https://doi.org/10.4171/owr/2024/41}.
\item Supplied literature assessment, checked on \textbf{2026-08-22}: the record labels the status open and describes an agenda concerning square-tiled surfaces and large-genus strata whose surrounding cited results do not settle the extracted program. This is supplied dated curation, not an independently verified literature assessment. The background expressly requests independent verification of time-sensitive status and cited claims.
\item Supplied statement-verification note: the extraction is labeled \emph{unrecoverable}. That note reports that the official report confirms a join between the end of one contribution and the title and opening sentence of the next, with no standalone problem; it says concrete square-tiled-surface questions occur in following records and that no clean statement is promoted. This draft preserves that supplied verification claim without claiming to have inspected the report.
\item Attribution retained from the extracted passage: Pinkham for the stated monomial-singularity and Weierstrass-semigroup relationship, and Kontsevich--Zorich for the referenced \(K(\pi,1)\) conjecture. No precise bibliographic references for either attribution, and no authoritative citation to a standalone problem for this ID, are supplied.
\end{itemize}
\end{document}
